Deux charges positives sont fixées sur une règle : $q_1 = +\qty{\v{q1}}{\micro C}$ à la position \qty{0}{cm} et $q_2 = +\qty{\v{q2}}{\micro C}$ à la position \qty{\v{d}}{cm}. On place une troisième charge $q_3 = +\qty{\v{q3}}{\micro C}$ entre les deux, à \qty{\v{x}}{cm} de $q_1$. La constante de Coulomb est $k = \qty{8.99e9}{N.m^2/C^2}$.

\begin{enumerate}[label=\alph*)]
    \item Déterminez la grandeur et l'orientation de la force exercée par $q_1$ sur $q_3$.
    \item Déterminez la grandeur et l'orientation de la force exercée par $q_2$ sur $q_3$.
    \item Déterminez la force résultante sur $q_3$ (grandeur et orientation).
\end{enumerate}
